\documentclass[titlepage,openany,a4paper,spanish]{book}

% Paquetes necesarios
\usepackage[utf8]{inputenc} 
\usepackage[T1]{fontenc}    
\usepackage[spanish,es-noshorthands]{babel}  
\usepackage{fancyhdr}         
\usepackage{amsmath, amssymb, amsfonts, amsthm}
\usepackage{hyperref}       
\usepackage{fullpage}       
\usepackage{titlesec}     
\usepackage{graphicx}       
\usepackage{float}   
\usepackage{ragged2e}          
\usepackage[backend=biber,style=numeric]{biblatex}
\usepackage{listings} % Para incluir código fuente
\usepackage{tikz}
\usepackage{pgfplots}
\usepackage{csquotes}
\pgfplotsset{compat=1.18}

\addbibresource{referencias.bib} 

% Configuración de encabezados y pies de página
\setlength{\headheight}{15.0pt}   % Ajuste de headheight para evitar advertencias

\fancypagestyle{plain}{
    \fancyhf{}
    \fancyhead[R]{        \begin{tikzpicture}[remember picture,overlay]
      \node[anchor=north east, xshift=-2.5cm, yshift=-1.5cm] at (current page.north east) {
          \includegraphics[width=3cm]{logo.png}
      };
  \end{tikzpicture}}
    \fancyhead[L]{TFG}
    \fancyfoot[C]{\thepage}
    \renewcommand{\headrulewidth}{0.2pt}
}

\pagestyle{plain}

% Redefinir el nombre de los capítulos de "Capítulo" a "Tema"
\renewcommand{\chaptername}{Tema}

% Personalizar el formato del capítulo para mostrar "Tema X: Título del Tema"
\titleformat{\chapter}[hang]
  {\bfseries\Huge}                   % Cambiado de \Huge a \LARGE para reducir tamaño
  {Capítulo \thechapter .}                 % Formato del número del capítulo
  {1em}                               % Espacio entre el número y el título
  {}                                  % Formato del título del capítulo

% Ajustar el espaciado antes y después del título del capítulo
\titlespacing*{\chapter}{0pt}{-1em}{1em} % Reduce el espacio antes del título

% Eliminación de márgenes adicionales
\setlength{\parindent}{0pt}          % Sin indentación en párrafos
\setlength{\headsep}{10pt}

% Comando personalizado (si es necesario)
\newcommand{\Lagr}{\mathcal{L}}

\begin{document}

% Inclusión de la portada

\begin{titlepage}
    \centering
    \vspace{1cm}
    
    % Logo de la universidad
    \includegraphics[height=5cm]{logo.png} 
    \vspace{1cm}
    
    % Nombre de la universidad
    {\LARGE \textbf{Universidad Europea de Madrid} \par}
    {\Large Escuela de Arquitectura, Ingeniería y Diseño \par}
    
    \vspace{3cm}
    
    % Línea decorativa superior
    \noindent\rule{12cm}{1.5pt}
    
    \vspace{0.8cm}
    
    % Título del trabajo
    {\huge \textbf{Trabajo Fin de Grado} \par}
    
    \vspace{0.5cm}
    
    {\Large Grado en Ingeniería Informática \par}
    
    \vspace{0.8cm}
    
    % Línea decorativa inferior
    \noindent\rule{12cm}{1.5pt}
    
    \vfill
    
    % Autor
    {\Large \textbf{Daniel Cuesta Sanz de Madrid} \par}
    
    \vspace{1cm}
    
    % Fecha
    {\large \today \par}
    
\end{titlepage}

\setcounter{page}{1}
\newpage


\textbf{\Large Título del Trabajo:} \Large TFG
\bigbreak
\textbf{\Large Autor:} {\Large Daniel Cuesta Sanz de Madrid}
\bigbreak
\textbf{\Large Tutor:} {\Large Nombre del tutor}
\bigbreak
\textbf{\Large Fecha:} \today 
\bigbreak

\normalsize 
\newpage

\section*{Resumen}
El presente Trabajo Fin de Grado\dots
\vspace{1cm}

\section*{Abstract}
This Final Degree Project aims to\dots

\newpage

\tableofcontents

\listoffigures

% Inclusión de Tema 1
\chapter{Introducción} % Define el título del capítulo
\section{Planteamiento del Problema}

\section{Objetivos del Proyecto}

\section{Estructura del proyecto}


\chapter{Objetivos}
\section{Objetivos generales}

\section{Objetivos específicos}

\newpage

\chapter{Referencias}
{$\space$\par}
\rhead[Referencias]{Referencias}

\defbibheading{myheading}[Referencias]{%
\chapter{Referencias}
}
\printbibliography[heading=myheading,title={Referencias}]


\end{document}
